% SEGMENT OLP-0090-S01
% Part: first-order-logic
% Chapter: natural-deduction
% Section: proving-things-quant

\documentclass[../../../include/open-logic-section]{subfiles}

\begin{document}

\olfileid{fol}{ntd}{prq}

\olsection{அளவையடைகளைக் கொண்ட \usetoken{P}{derivation}}

\begin{ex}
அளவையடைகளைக் கையாளும்போது தனிமாறி நிபந்தனையை மீறாமல் இருப்பதை உறுதிசெய்ய
வேண்டும். சில அனுமானங்களைச் செயல்படுத்தும் வரிசையை மாற்றிப் பார்க்கவும்
சிலவேளைகளில் இது நம்மைக் கோருகிறது. பொதுவாக, தனிமாறி நிபந்தனைக்கு உட்பட்ட
விதிகளை முதலில் கவனிக்க முயல்வது உதவும் (நிறைவுற்ற நிறுவலில் அவை கீழ்ப்பகுதியில்
இருக்கும்).

% SEGMENT OLP-0090-S02
!!{formula}~$\lexists[x][\lnot !A(x)] \lif \lnot \lforall[x][!A(x)]$ இன்
!!a{derivation} ஐ எவ்வாறு கொடுப்போம் எனப் பார்ப்போம். வழக்கம்போல் தொடங்கி,
பின்வருமாறு எழுதுகிறோம்:
\begin{prooftree}
\AxiomC{}
\UnaryInfC{$\lexists[x][\lnot !A(x)]\lif \lnot \lforall[x][!A(x)]$}
\end{prooftree}
அந்தக் கடைசிப் படியை \Intro{\lif} விதியால் நியாயப்படுத்த என்ன தேவைப்படும்
என்பதை முதலில் எழுதுகிறோம்.
\begin{prooftree}
\AxiomC{$\Discharge{\lexists[x][\lnot !A(x)]}{1}$}
\DeduceC{$\lnot \lforall[x][!A(x)]$}
\DischargeRule{\Intro{\lif}}{1}
\UnaryInfC{$\lexists[x][\lnot !A(x)]\lif \lnot \lforall[x][!A(x)]$}
\end{prooftree}
$\lnot \lforall[x][!A(x)]$ க்கு நேரடியாகப் பயன்படுத்தத்தக்க விதி தெளிவாகத்
தெரியாததால், \Elim{\lexists} விதியைப் பயன்படுத்தும் வகையில் !!{derivation} ஐ
அமைப்போம். இங்கு தனிமாறி நிபந்தனையைக் கவனித்து,
$\lexists[x][!A(x)]$ இலோ அது சார்ந்திருக்கும் எந்த எடுகோளிலோ இடம்பெறாத ஒரு
மாறிலிக் குறியைத் தேர்ந்தெடுக்க வேண்டும். (ஆனால் !!{constant}s எவையும்
இடம்பெறாததால் எந்தத் தேர்வும் பொருந்தும்.)
\begin{prooftree}
\AxiomC{$\Discharge{\lexists[x][\lnot !A(x)]}{1}$}
\AxiomC{$\Discharge{\lnot !A(a)}{2}$}
\DeduceC{$\lnot \lforall[x][!A(x)]$}
\DischargeRule{\Elim{\lexists}}{2}
\BinaryInfC{$\lnot \lforall[x][!A(x)]$}
\DischargeRule{\Intro{\lif}}{1}
\UnaryInfC{$\lexists[x][\lnot !A(x)] \lif \lnot \lforall[x][!A(x)]$}
\end{prooftree}

% SEGMENT OLP-0090-S03
$\lnot \lforall[x][!A(x)]$ ஐ வருவிக்க \Intro{\lnot} விதியைப் பயன்படுத்த
முயல்வோம். கூடுதல் எடுகோளாக $\lforall[x][!A(x)]$ ஐப் பயன்படுத்தியும் இருக்கக்
கூடிய ஒரு முரண்பாட்டை நாம் வருவிக்க வேண்டும். இந்த முரண்பாடு,
\Elim{\lexists} அனுமானத்தால் விடுவிக்கப்படும் எடுகோள்~$\lnot !A(a)$ ஐயும்
உள்ளடக்கலாம். பின்வருமாறு அமைக்கலாம்:
\begin{prooftree}
\AxiomC{$\Discharge{\lexists[x][\lnot !A(x)]}{1}$}
\AxiomC{$\Discharge{\lnot !A(a)}{2}, \Discharge{\lforall[x][!A(x)]}{3}$}
\DeduceC{$\lfalse$}
\DischargeRule{\Intro{\lnot}}{3}
\UnaryInfC{$\lnot \lforall[x][!A(x)]$}
\DischargeRule{\Elim{\lexists}}{2}
\BinaryInfC{$\lnot \lforall[x][!A(x)]$}
\DischargeRule{\Intro{\lif}}{1}
\UnaryInfC{$\lexists[x][\lnot !A(x)]\lif \lnot \lforall[x][!A(x)]$}
\end{prooftree}

% SEGMENT OLP-0090-S04
முரண்பாட்டைப் பெறுவதற்கு நெருங்கிவிட்டதாகத் தெரிகிறது. பயன்படுத்துவதற்கு
மிக எளிய விதி தனிமாறி நிபந்தனை ஏதும் இல்லாத \Elim{\lforall}. அனைத்தளவிட்ட
$x$ க்குப் பதிலாக நாம் விரும்பும் எந்தச் சொல்லையும் பயன்படுத்தலாம்; ஆகவே
முரண்பாட்டை அடையும் வகையில்~$a$ ஐயே தொடர்ந்து பயன்படுத்துவது பொருத்தமானது.
\begin{prooftree}
\AxiomC{$\Discharge{\lexists[x][\lnot !A(x)]}{1}$}
\AxiomC{$\Discharge{\lnot !A(a)}{2}$}
\AxiomC{$\Discharge{\lforall[x][!A(x)]}{3}$}
\RightLabel{\Elim{\lforall}}
\UnaryInfC{$!A(a)$}
\RightLabel{\Elim{\lnot}}
\BinaryInfC{$\lfalse$}
\DischargeRule{\Intro{\lnot}}{3}
\UnaryInfC{$\lnot \lforall[x][!A(x)]$}
\DischargeRule{\Elim{\lexists}}{2}
\BinaryInfC{$\lnot \lforall[x][!A(x)]$}
\DischargeRule{\Intro{\lif}}{1}
\UnaryInfC{$\lexists[x][\lnot !A(x)]\lif \lnot \lforall[x][!A(x)]$}
\end{prooftree}

% SEGMENT OLP-0090-S05
குறிப்பாக அளவையடைகளைக் கையாளும்போது, தனிமாறி நிபந்தனை மீறப்படவில்லை என்பதை
இந்த நிலையில் மீண்டும் சரிபார்ப்பது முக்கியம். தனிமாறி நிபந்தனைக்கு உட்பட்டு
நாம் பயன்படுத்திய ஒரே விதி \Elim{\lexists}; அதன் தனிமாறி~$a$ அது சார்ந்த எந்த
% SOURCE CORRECTION TA-ND-004: raw \exists is corrected to \lexists.
எடுகோளிலும் இடம்பெறவில்லை. ஆகவே இது சரியான !!{derivation}.
\end{ex}

% SEGMENT OLP-0090-S06
\begin{ex}
சிலவேளைகளில் சில !!{formula}s இலிருந்து மற்றோர் !!a{formula} ஐ வருவிக்கலாம்.
இந்நிலைகளில் விடுவிக்கப்படாத எடுகோள்கள் இருக்கலாம். நமது இறுதி இலக்குடன்
சேர்த்து எடுகோள்களையும் கண்காணிப்பது முக்கியம்.

எடுகோள்கள்~$\lexists[x][(!A(x) \land !B(x))]$ மற்றும்
$\lforall[x][(!B(x) \lif !C(x,b))]$ ஆகியவற்றிலிருந்து
!!{formula}~$\lexists[x][!C(x,b)]$ இன் !!a{derivation} ஐ எவ்வாறு கொடுப்போம்
எனப் பார்ப்போம். வழக்கம்போல் முடிவை அடியில் எழுதுகிறோம்.
\begin{prooftree}
\AxiomC{}
\UnaryInfC{$\lexists[x][!C(x,b)]$}
\end{prooftree}

% SEGMENT OLP-0090-S07
நாம் பயன்படுத்துவதற்கு இரண்டு முன்கோள்கள் உள்ளன. முதலாவதைப் பயன்படுத்த,
அதாவது $\lexists[x][(!A(x) \land !B(x))]$ இலிருந்து
$\lexists[x][!C(x, b)]$ இன் !!a{derivation} ஐக் கண்டறிய முயல, \Elim{\lexists}
விதியைப் பயன்படுத்துவோம். அதற்குத் தனிமாறி நிபந்தனை இருப்பதால் அவ்விதியை
முதலில் பயன்படுத்துகிறோம். பின்வருவது கிடைக்கிறது:
\begin{prooftree}
\AxiomC{$\lexists[x][(!A(x) \land !B(x))]$}
\AxiomC{$\Discharge{!A(a) \land !B(a)}{1}$}
\DeduceC{$\lexists[x][!C(x,b)]$}
\DischargeRule{\Elim{\lexists}}{1}
\BinaryInfC{$\lexists[x][!C(x,b)]$}
\end{prooftree}

% SEGMENT OLP-0090-S08
நாம் கையாளும் இரு எடுகோள்களிலும்~$!B$ பொதுவாக உள்ளது. இப்போது
\Elim{\land} ஐப் பயன்படுத்தி~$!B(a)$ ஐத் தனியாகப் பிரித்தெடுப்பது பயனுள்ளதாக
இருக்கலாம்.
\begin{prooftree}
\AxiomC{$\lexists[x][(!A(x) \land !B(x)])$}
\AxiomC{$\Discharge{!A(a) 
\land !B(a)}{1}$}
\RightLabel{\Elim{\land}}
\UnaryInfC{$!B(a)$}
\DeduceC{$\lexists[x][!C(x,b)]$}
\DischargeRule{\Elim{\lexists}}{1}
\BinaryInfC{$\lexists[x][!C(x,b)]$}
\end{prooftree}

% SEGMENT OLP-0090-S09
நாம் பயன்படுத்த வேண்டிய இரண்டாம் எடுகோள்
$\lforall[x][(!B(x) \lif !C(x,b))]$. தனிமாறி நிபந்தனை இல்லாததால்,
\Elim{\lforall} ஐப் பயன்படுத்தி $x$ க்கு !!{constant}~$a$ ஐ இட்டு
$!B(a) \lif !C(a, b)$ ஐப் பெறலாம். இப்போது $!B(a) \lif !C(a,b)$,
$!B(a)$ இரண்டும் நம்மிடம் உள்ளன. அடுத்த படியாக \Elim{\lif} விதியை
நேரடியாகப் பயன்படுத்த வேண்டும்.
\begin{prooftree}
\AxiomC{$\lexists[x][(!A(x) \land !B(x))]$}
\AxiomC{$\lforall[x][(!B(x) \lif !C(x,b))]$}
\RightLabel{\Elim{\lforall}}
\UnaryInfC{$!B(a) \lif !C(a,b)$}
\AxiomC{$\Discharge{!A(a) 
\land !B(a)}{1}$}
\RightLabel{\Elim{\land}}
\UnaryInfC{$!B(a)$}
\RightLabel{\Elim{\lif}}
\BinaryInfC{$!C(a,b)$}
\DeduceC{$\lexists[x][!C(x,b)]$}
\DischargeRule{\Elim{\lexists}}{1}
\insertBetweenHyps{\hspace{-5em}}
\BinaryInfC{$\lexists[x][!C(x,b)]$}
\end{prooftree}

% SEGMENT OLP-0090-S10
இலக்கிற்கு மிக அருகில் வந்துவிட்டோம்!{} \Intro{\lexists} இன் ஒரே பயன்பாட்டில்
இலக்கை அடைகிறோம்.
\begin{prooftree}
\AxiomC{$\lexists[x][(!A(x) \land !B(x))]$}
\AxiomC{$\lforall[x][(!B(x) \lif !C(x,b))]$}
\RightLabel{\Elim{\lforall}}
\UnaryInfC{$!B(a) \lif !C(a,b)$}
\AxiomC{$\Discharge{!A(a) 
\land !B(a)}{1}$}
\RightLabel{\Elim{\land}}
\UnaryInfC{$!B(a)$}
\RightLabel{\Elim{\lif}}
\BinaryInfC{$!C(a,b)$}
\RightLabel{\Intro{\lexists}}
\UnaryInfC{$\lexists[x][!C(x,b)]$}
\DischargeRule{\Elim{\lexists}}{1}
\insertBetweenHyps{\hspace{-5em}}
\BinaryInfC{$\lexists[x][!C(x,b)]$}
\end{prooftree}
ஒவ்வொரு படியிலும் தனிமாறி நிபந்தனைகள் மீறப்படவில்லை என்பதை உறுதிசெய்ததால்,
இது சரியான !!{derivation} என நம்பலாம்.
\end{ex}

% SEGMENT OLP-0090-S11
\begin{ex}
எடுகோள்கள்~$\lforall[x][!A(x)] \lif \lexists[y][!B(y)]$ மற்றும்
$\lnot\lexists[y][!B(y)]$ ஆகியவற்றிலிருந்து
!!{formula}~$\lnot\lforall[x][!A(x)]$ இன் !!a{derivation} ஐக் கொடுக்கவும்.
வழக்கம்போல் இலக்கு !!{formula} வை அடியில் எழுதுகிறோம்.
\begin{prooftree}
\AxiomC{}
\UnaryInfC{$\lnot\lforall[x][!A(x)]$}
\end{prooftree}

% SEGMENT OLP-0090-S12
!!{derivation}[gen] கடைசி வரி ஒரு மறுப்பு; ஆகவே \Intro{\lnot} ஐப் பயன்படுத்த
முயல்வோம். இதற்கு ஒரு முரண்பாட்டை எவ்வாறு !!{derive} செய்வது என்று கண்டறிய
வேண்டும்.
\begin{prooftree}
\AxiomC{$\Discharge{\lforall[x][!A(x)]}{1}$}
\DeduceC{$\lfalse$}
\DischargeRule{\Intro{\lnot}}{1}
\UnaryInfC{$\lnot\lforall[x][!A(x)]$}
\end{prooftree}

% SEGMENT OLP-0090-S13
இதுவரை சரியாக உள்ளது. \Elim{\lforall} ஐப் பயன்படுத்தலாம்; ஆனால் இலக்கை
அடைய அது உதவுமா என்பது தெளிவில்லை. அதற்குப் பதிலாக நமது எடுகோள்களில் ஒன்றைப்
பயன்படுத்துவோம். $\lforall[x][!A(x)] \lif \lexists[y][!B(y)]$ உடன்
$\lforall[x][!A(x)]$ ஐச் சேர்த்தால் \Elim{\lif} விதியைப் பயன்படுத்தலாம்.
\begin{prooftree}
\AxiomC{$\lforall[x][!A(x)] \lif \lexists[y][!B(y)]$}
\AxiomC{$\Discharge{\lforall[x][!A(x)]}{1}$}
\RightLabel{\Elim{\lif}}
\BinaryInfC{$\lexists[y][!B(y)]$}
\DeduceC{$\lfalse$}
\DischargeRule{\Intro{\lnot}}{1}
\UnaryInfC{$\lnot\lforall[x][!A(x)]$}
\end{prooftree}

% SEGMENT OLP-0090-S14
இப்போது பயன்படுத்துவதற்கு ஓர் இறுதி எடுகோள் உள்ளது; \Elim{\lnot} ஐப்
பயன்படுத்தி முரண்பாட்டை அடைய அது உதவும் எனத் தெரிகிறது.
\begin{prooftree}
\AxiomC{$\lnot\lexists[y][!B(y)]$}
\AxiomC{$\lforall[x][!A(x)] \lif \lexists[y][!B(y)]$}
\AxiomC{$\Discharge{\lforall[x][!A(x)]}{1}$}
\RightLabel{\Elim{\lif}}
\BinaryInfC{$\lexists[y][!B(y)]$}
\RightLabel{\Elim{\lnot}}
\BinaryInfC{$\lfalse$}
\DischargeRule{\Intro{\lnot}}{1}
\UnaryInfC{$\lnot\lforall[x][!A(x)]$}
\end{prooftree}
\end{ex}

% SEGMENT OLP-0090-S15
\begin{prob}
பின்வருவனவற்றைக் காட்டும் !!{derivation}s ஐக் கொடுக்கவும்:
\begin{enumerate}
\item $\Proves (\lforall[x][!A(x)] \land \lforall[y][!B(y)]) \lif
\lforall[z][(!A(z) \land !B(z))]$.
\item $\Proves (\lexists[x][!A(x)] \lor \lexists[y][!B(y)]) \lif
\lexists[z][(!A(z) \lor !B(z))]$.
\item $\lforall[x][(!A(x) \lif !B)] \Proves \lexists[y][!A(y)] \lif !B$.
\item $\lforall[x][\lnot !A(x)] \Proves \lnot\lexists[x][!A(x)]$.
\item $\Proves \lnot\lexists[x][!A(x)] \lif \lforall[x][\lnot !A(x)]$.
\item $\Proves \lnot\lexists[x][\lforall[y][((!A(x,y) \lif \lnot
!A(y,y)) \land (\lnot !A(y,y) \lif !A(x,y)))]]$.
\end{enumerate}
\end{prob}

% SEGMENT OLP-0090-S16
\begin{prob}
பின்வருவனவற்றைக் காட்டும் !!{derivation}s ஐக் கொடுக்கவும்:
\begin{enumerate}
\item $\Proves \lnot\lforall[x][!A(x)] \lif \lexists[x][\lnot!A(x)]$.
\item $(\lforall[x][!A(x)] \lif !B) \Proves \lexists[y][(!A(y) \lif !B)]$.
\item $\Proves \lexists[x][(!A(x) \lif \lforall[y][!A(y)])]$.
\end{enumerate}
(இவை அனைத்திற்கும் $\FalseCl$~விதி தேவை.)
\end{prob}

\end{document}
